\section{Virtual threads}
\label{sec:virtual-threads}

O \citetitle{doc:jep444} tornou as virtual threads finais no JDK~21. Uma
virtual thread é uma \texttt{java.lang.Thread} que não fica presa a uma
thread do sistema operacional: quando ela bloqueia em I/O, a JVM a
\termo{desmonta}{\emph{Unmount}: a virtual thread sai da
\emph{carrier thread} (a thread de plataforma que a executa) e esta fica
livre para rodar outra virtual thread.} e usa a mesma thread de
plataforma para outra tarefa.

\lstinputlisting[language=Java]{generated/source/ThreadsVirtuais.java}

\inputbytecode{ThreadsVirtuais}

\paragraph{Lendo a listagem.}
Não há nenhuma instrução especial. O offset 0 é uma chamada estática
comum a \texttt{Executors.newVirtualThreadPerTaskExecutor}, o offset 17
é o \texttt{invokedynamic} de um lambda (que captura \texttt{tarefa},
por isso o descritor \nolinkurl{(I)Ljava/util/concurrent/Callable;}) e os
offsets 34--68 são o \texttt{try}-with-resources que fecha o executor (o
\texttt{close()} espera as tarefas). Toda a diferença entre virtual e
plataforma está no runtime do JDK, não no seu bytecode: o código da
tarefa é o mesmo; o que muda é quem cria a thread que o executa.

\paragraph{As decisões que isso muda.}
A documentação da Oracle (\citetitle{doc:oracle-virtual-threads}) e o
JEP~444 dão regras diretas:
\begin{itemize}
\item Virtual threads servem para \emph{escala} (muitas tarefas que passam
  a maior parte do tempo esperando), não para velocidade: \enquote{Virtual
  threads are not faster threads}. Para trabalho de CPU longo, não ajudam.
\item Não faça pool de virtual threads. Uma por tarefa; para limitar
  concorrência a um recurso escasso (conexões com um serviço), use um
  \texttt{Semaphore}, não um pool de tamanho fixo.
\item Cuidado com \texttt{ThreadLocal} usado como cache de objetos caros:
  cada virtual thread roda uma só tarefa, então o cache nunca é
  reaproveitado. Para passar contexto imutável, o
  \citetitle{doc:jep506} (final no JDK~25) oferece \texttt{ScopedValue}.
\end{itemize}

\paragraph{Pinning e \texttt{synchronized}.}
No JDK~21, uma virtual thread que bloqueia dentro de um bloco
\texttt{synchronized} fica \termo{pinned}{Presa à carrier thread: não
consegue desmontar enquanto bloqueia, e a carrier fica ocupada.}. O
JEP~444 diz que isso não torna o programa incorreto, mas pode limitar a
escala, e recomendava trocar \texttt{synchronized} por
\texttt{ReentrantLock} em trechos frequentes que guardam I/O. O
\citetitle{doc:jep491} (JDK~24) removeu quase todo esse pinning: a
virtual thread agora desmonta mesmo dentro de \texttt{synchronized}. Ainda
fica presa quando código nativo chama de volta Java e bloqueia, e o
evento JFR \texttt{jdk.VirtualThreadPinned} continua existindo para esses
casos. A seção~\ref{sec:sincronizacao} mostra o
\texttt{monitorenter} por trás disso.

Consequência prática: se o seu serviço roda em JDK~21, o conselho de
trocar \texttt{synchronized} vale; em JDK~24 ou mais novo, não é mais
motivo para refatorar.

\paragraph{Ponte com Quarkus.}
O guia \citetitle{doc:quarkus-virtual-threads} mostra a anotação
\texttt{@RunOnVirtualThread}, que faz o Quarkus invocar o método anotado
numa nova virtual thread, e exige Java~21 ou mais novo. O guia repete os
pontos acima no contexto do Quarkus: em Java~24+, \texttt{synchronized}
não causa mais pinning; uma computação longa monopoliza a carrier até
terminar; e pool de objetos via \texttt{ThreadLocal} é contraproducente
com virtual threads. A Parte~III volta a esse assunto.
